No single community signal captures all dimensions of community of interest
recognized in redistricting law.
Economic character (M.1), physical land use (M.2), housing character (M.3),
commuting shed membership (M.4), topographic community (M.5), administrative
zone co-membership (M.6), and transit accessibility (M.7) each measure a
distinct dimension of community bond.
Plans drawn with any single signal are defensible on that dimension and
potentially vulnerable on others.
Plans drawn with the composite index are defensible on all dimensions
simultaneously.

This paper presents M.8, the concluding paper in the M-track community
character series, which defines the Composite Community Character Index:
a weighted average of all available M.1--M.7 similarity signals, designed
for use in redistricting proceedings as a court-admissible, quantitative
operationalization of communities of interest.

The composite formula weights zone co-membership (M.6) most highly at 0.30,
reflecting its status as the most legally recognized community signal
(school district and municipal service co-membership is explicitly cited
in redistricting criteria in California, Colorado, and Arizona).
Economic character (M.1) and land use (M.2) share 0.20 each, housing character
(M.3) receives 0.15, commuting shed (M.4) receives 0.10, and topographic
and transit signals (M.5, M.7) receive 0.03 and 0.02 respectively.
The formula normalizes over available signals: missing data for a component
causes graceful degradation rather than fallback to a uniform default.

The court usage template in Section~4 provides a complete expert witness
statement with specific language for redistricting briefs and testimony.
The template addresses data provenance, criterion alignment, community
identification, algorithmic transparency, and the statement that no
electoral, racial, or partisan data entered any component signal.
All empirical results are deferred to Phase~2 pending component data downloads.
